% Source pp. 173-181 (Lectures13-18.pdf).
\section{More on Wilson's RG}

The last time we have defined the \emph{Wilson's Renormalization Group
Transformation}. This transformation acts on the action functionals
$\mathcal A[\varphi_0]$,
\[
\mathcal A\to\mathcal A'\,,
\]
which represent \emph{cutoff field theories}, i.e.\ field theories equipped with
some cutoff regularization, with cutoff parameter $\Lambda$. The transformation
depends on a parameter $L\ge1$ and a function $Z(L)$, and consists of two steps.

\paragraph{STEP 1:} Integrating out the ``fast'' degrees of freedom.

The field $\varphi_0(x)$ is split into its ``slow'' and ``fast'' parts,
\[
\varphi_0(x)=\varphi_1(x)+\tilde\varphi(x)\,,
\]
where
\begin{equation}\label{eq:slowfast}
\varphi_1(x)=\int_{|k|<\Lambda/L}\phi(k)\,e^{ikx}\,,\qquad
\tilde\varphi(x)=\int_{\Lambda/L<|k|<\Lambda}\phi(k)\,e^{ikx}\,.
\end{equation}
Integrating the part $\tilde\varphi$ out one defines new action
\begin{equation}\label{eq:A1def}
\Const\,e^{-\mathcal A_1[\varphi_1]}=\int[D\tilde\varphi]\,e^{-\mathcal A[\varphi_1+\tilde\varphi]}
\end{equation}
which depends only on the ``slow'' degrees of freedom, with the cutoff parameter
$\Lambda/L$.

\paragraph{STEP 2:} Rescaling of the coordinates $x\to x/L$.

One makes the field variable transformation
\begin{equation}\label{eq:rescale}
\varphi_1(x)=Z^{-\frac12}(L)\,\varphi_0(x/L)
\end{equation}
in the functional integral over $\varphi_1$, defining new integration variable
$\varphi_0$. This transformation returns the cutoff parameter to its original
value $\Lambda$, and yields new action
\begin{equation}\label{eq:Aprime}
\mathcal A'[\varphi_0(x)]=\mathcal A_1\bigl[Z^{-\frac12}(L)\,\varphi_0(x/L)\bigr]\,.
\end{equation}

In principle, the factor $Z^{\frac12}(L)$ in \eqref{eq:rescale}, \eqref{eq:Aprime}
is arbitrary. Field renormalizations $\varphi_0\to Z^{\frac12}\varphi_0$ do not
change physical content of field theory, it can be thought of as some sort of
gauge transformation.\footnote{Much more general class of ``gauge
transformations'' can be applied to field theories, the so called ``field
redefinitions''. One considers generic local change of variables of the form
$\varphi_0(x)\to f\bigl(\varphi_0(x),\partial_\mu\varphi_0(x),\cdots\bigr)$, which
leaves the quasilocal nature of the action \eqref{eq:quasilocal} intact. Here we
will not consider implications of this huge ambiguity. Of course, it is not to be
mixed with the usual gauge transformations in gauge field theories.} When we
characterize field theories by their action functionals, two actions different
only in such renormalization are equivalent,
\[
\mathcal A\bigl[Z^{\frac12}\varphi_0\bigr]\sim\mathcal A[\varphi_0]\,.
\]
In practice this redundancy is usually eliminated by some field normalization
condition which plays role of gauge fixing condition. In this case the factor
$Z^{\frac12}(L)$ must be adjusted so that the transformation
$\mathcal A\to\mathcal A'$ does preserve that condition.

It follows from the above definition that if $\mathcal A'$ is obtained from
$\mathcal A$ by the Wilson's RG transformation, the corresponding correlation
functions are related as,
\[
\text{f.t.}\,\bigl\langle\varphi_0(x_1)\cdots\varphi_0(x_n)\bigr\rangle_{\mathcal A'}
=\text{f.t.}\,Z^{\frac n2}(L)\,\bigl\langle\varphi_0(x_1L)\cdots\varphi_0(x_nL)\bigr\rangle_{\mathcal A}
\]
where it is understood that the Fourier transforms of the correlation functions
are taken and the comparison is made at low values of the momenta
$|p_i|\ll\Lambda$. In other words, the theory described by the action
$\mathcal A'$ differs from that described by $\mathcal A$ by a scale
transformation $x\to x/L$ only. The Wilson's RG shows what it takes to make such
scale transformation in field theory. Just rescaling the coordinates in the
action, as is done in \eqref{eq:rescale}, is not sufficient. The ``naive''
scaling transformation, represented by the step 2 above, must be supplemented by
the integration over the fast modes in the step 1. Roughly speaking, the dilation
$x\to x/L$ increases the spatial density of the field degrees of freedom, and one
has to ``dilute'' them back to the original density by performing the step 1. In
this sense the step 1 can be thought of as representing the transformation of the
measure in the functional integral.

One remark is in order here. In general, implementation of the step 1 of the
Wilson's RG (elimination of the fast modes) depends on particular form of the
cutoff regularization implied. The definition \eqref{eq:slowfast} corresponds to
the ``sharp momentum cutoff''. We know that other cutoff regularizations are
possible. For example, in the lattice regularization, where the field degrees of
freedom are associated with the lattice cells, the procedure described above is
hardly natural. Instead, one can consider a sublattice whose cells consist of
$L^d$ elementary cells ($L$ must be integer in this case). The ``slow'' field
$\varphi_1$ is then defined as the original field averaged over these greater
cells, and the step 1 amounts to integrating over $\varphi_0$ with $\varphi_1$
fixed. I am mentioning this here to stress that the two step procedure of the
Wilson's RG can be considered as a general idea rather than an exact
prescription, the idea whose implementation depends on the problem in hand.
However, in this discussion we will follow the particular implementation
described in \eqref{eq:slowfast}.

We have seen that even if the original action $\mathcal A$ has some simple form,
say $\varphi^4$, the result of the Wilson's RG transformation is usually a
complicated ``quasilocal'' action which contains all possible terms with
arbitrarily high powers of the field and with arbitrarily high powers and orders
of the derivatives. On the other hand, if $\mathcal A$ is generic quasilocal
action, one can expect that the action $\mathcal A'$ will again be of the
quasilocal form. Under this assumption, the Wilson's RG transformation acts in
the space of all quasilocal action functionals. In what follows I will use the
notation $\Sigma$ for this infinite-dimensional space, so that
\[
RG:\ \Sigma\to\Sigma\,.
\]
Also, I will denote the RG transformation as
\[
\mathcal A'=RG_l\{\mathcal A\}\,,
\]
using more convenient parameter
\[
l=\log L\ge0
\]
instead of $L$. It is not difficult to understand that these transformations
obey the composition law
\begin{equation}\label{eq:semigroup}
RG_{l+l'}\{\mathcal A\}=RG_l\{RG_{l'}\{\mathcal A\}\}\,,\qquad RG_0\{\mathcal A\}=\mathcal A\,.
\end{equation}
These properties define one-parameter \emph{semigroup}. It is not a group
because a group contains inverse elements, and inverse to $RG_l$ in general
cannot be defined. The reason is simple but important: the step 1 of the RG
transformation leads to a loss of information about the interactions at the
cutoff scales; this step cannot be ``undone''.

In general exact evaluation of the RG transformation for arbitrary $l$ is at
least as difficult as exact evaluation of the full functional integral, i.e.\
usually hopeless. If the interaction is not small no reliable approximations are
even available. One can hope however to find good approximations for $RG_l$ with
$l$ small enough. Then according to the composition law \eqref{eq:semigroup} the
transformation $RG_{Nl}$ can be obtained as $N$th iteration of $RG_l$. While it
is not clear how to evaluate the functional integral at once, maybe it is
possible to do it step by step, integrating out the shortest scale degrees of
freedom, and then repeating this again and again.

It is usually most convenient to consider the infinitesimal RG transformation
\[
RG_{\delta l}\{\mathcal A\}=\mathcal A+B\{\mathcal A\}\,\delta l+O(\delta l^2)\,,
\qquad\text{where}\qquad
B\{\mathcal A\}=\frac{\dd}{\dd l}\,RG_l\{\mathcal A\}\Big|_{l=0}\,.
\]
It follows from \eqref{eq:semigroup} that the one-parameter family of actions
\[
\mathcal A_l=RG_l\{\mathcal A_0\}\,,
\]
where $\mathcal A_0$ is an arbitrary ``initial'' action, satisfies the differential
equation
\begin{equation}\label{eq:WilsonFlow}
\frac{\dd}{\dd l}\,\mathcal A_l=B\{\mathcal A_l\}\,,
\end{equation}
which is usually referred to as \emph{Wilson's RG flow equation}. It generates the
scale transformations in the space of quasilocal actions $\Sigma$, namely, two
actions $\mathcal A_{l_1}$ and $\mathcal A_{l_2}$ belonging to the same integral
curve $\mathcal A_l$ (sometimes called the \emph{RG trajectory}) describe the
theories which differ by the scale transform $x\to x\exp(l_1-l_2)$. Thus the
problem splits into two parts: (i) Evaluation of the RG operator
$B\{\mathcal A\}$ and (ii) Integration of the RG flow equation
\eqref{eq:WilsonFlow}.

The RG flow equation \eqref{eq:WilsonFlow} is written above in very abstract
form. To make it substantial one has to introduce some coordinates on $\Sigma$.
Generic quasilocal action can be written as (I assume $\varphi_0\to-\varphi_0$
symmetry)
\begin{multline}\label{eq:quasilocal}
\mathcal A[\varphi_0]=\int \dd^dx\,\Bigl[\sum_{k=1}^\infty\frac{u_{2k}}{2k!}\,\varphi_0^{2k}(x)+\sum_{k=0}^\infty\frac{v_{2k}}{2k!}\,\varphi_0^{2k}\,(\partial\varphi_0)^2(x)\\
+\sum_{k=0}^\infty\frac{w_{2k+1}}{2k!}\,\varphi_0^{2k+1}\,(\partial\varphi_0)^2\,\partial^2\varphi_0(x)+\cdots\Bigr]\,.
\end{multline}
The coefficients (``coupling constants'') $u_{2k},v_{2k},w_{2k+1},\dots$ can be
taken as such coordinates. In general there are infinitely many couplings and
therefore the space $\Sigma$ is infinite-dimensional. One of these parameters
can be eliminated by appropriate renormalization of $\varphi_0$, for instance one
can use this freedom to set
\[
v_0=1
\]
(the ``gauge fixing'' discussed above).

In terms of these coordinates the RG flow equation \eqref{eq:WilsonFlow} can be
written as
\begin{equation}\label{eq:Bfunctions}
\begin{aligned}
\frac{\dd}{\dd l}\,u_{2k}(l)&=B_{u_{2k}}\bigl(\{u_{2k}(l)\},\{v_{2k}(l)\},\{w_{2k+1}(l)\},\dots\bigr)\,;\\
\frac{\dd}{\dd l}\,v_{2k}(l)&=B_{v_{2k}}\bigl(\{u_{2k}(l)\},\{v_{2k}(l)\},\{w_{2k+1}(l)\},\dots\bigr)\,;\\
&\ \ \text{etc}
\end{aligned}
\end{equation}
where $u_{2k}(l),v_{2k}(l),\dots$ are the coupling constants corresponding to the
action $\mathcal A_l$, and $B$'s are some functions of all these infinitely many
coordinates. Of course it is impossible to calculate exactly all these functions.
Neither can one hope to integrate exactly all the RG flow equations. In practical
calculations some approximations are necessary. Most importantly one must find
some finite-dimensional approximation to the infinite-dimensional space $\Sigma$.
But before doing any calculations it is useful to think about what are the
properties of the RG flow we must be interested in.

\subsection*{Topological properties of RG flow\protect\footnote{The discussion in
this section follows closely that in K.~Wilson, J.~Kogut, ``The RG and epsilon
expansion'', Phys.\ Rep.\ \textbf{12C} (1974).}}

Some general topological properties of RG flow can be discussed in an abstract
form without reference to specific form of the RG transformation. Some general
assumptions will be made concerning the actions generated by the RG
transformations. Most importantly, we will make an assumption already mentioned
above, namely, that the RG transformation $RG_l$ applied to a quasilocal action
$\mathcal A$ produces an action $\mathcal A_l$ which is quasilocal. In other
words, our main assumption is that $RG:\Sigma\to\Sigma$. This is not a trivial
assumption. I remind you that our term ``quasilocal action'' means not only an
action which admits formal derivative expansion \eqref{eq:quasilocal}, but also
that this expansion converges for all fields $\varphi_0$ which contain only
Fourier components with $|k|\lesssim\Lambda$. Roughly speaking quasilocal action
can contain nonlocal interactions of the range $\sim\Lambda^{-1}$ but not much
larger. There is no proof to this statement, but it holds for many known
implementations of the RG procedure, so we accept it.

Under this assumption, if $\mathcal A_0\in\Sigma$, the associated RG trajectory
$\mathcal A_l$, i.e.\ the solution of the flow equation \eqref{eq:WilsonFlow}
with the initial condition $\mathcal A_{l=0}=\mathcal A_0$, lies inside $\Sigma$
for all $l>0$ (figure~\ref{fig:trajectory}).
\begin{figure}[ht]
\centering\figfile{trajectory}
\caption{An RG trajectory $\mathcal A_l$ starting at $\mathcal A_0$ stays inside
the space $\Sigma$ of quasilocal actions.}
\label{fig:trajectory}
\end{figure}

The trajectory $\mathcal A_l$, as a function of the parameter $l$, gives a set of
actions which all describe essentially the same physical system, the only
difference being in the coordinate rescalings. Namely one must make an $l$
dependent change of scale when relating $\mathcal A_l$ to physics.

It is useful to discuss in terms of physical mass $m$, or, as it is more common in
statistical mechanics, in terms of the ``correlation length'' $R_c=m^{-1}$. This
is defined as follows. Consider the two-point correlation function. Its typical
large-distance behavior in Euclidean space is (recall for instance your solution
of the behavior of the Klein-Gordon propagator $D(x)$ at $|x|\to\infty$)
\[
\bigl\langle\varphi_0(x)\varphi_0(0)\bigr\rangle\simeq e^{-\frac{|x|}{R_c}}\qquad\text{as}\qquad|x|\to\infty\,.
\]
Let $R_c$ be the correlation length associated with our system described by the
action $\mathcal A_0$, with the cutoff $\Lambda$. Of course $\mathcal A_l$
describes the same system, but in the units in which the cutoff associated with
$\mathcal A_l$ is still $\Lambda$, the correlation length will be
\[
R_c(l)=e^{-l}\,R_c\,.
\]
One can define various subspaces, or ``surfaces'', in $\Sigma$. The elements of
$\Sigma$ can be separated according to their correlation length.
Figure~\ref{fig:Rcsurfaces} shows schematically the surfaces of constant $R_c$.
\begin{figure}[ht]
\centering\figfile{Rcsurfaces}
\caption{Surfaces of constant correlation length in $\Sigma$. The thick curve is
the critical surface $R_c=\infty$.}
\label{fig:Rcsurfaces}
\end{figure}

Of special interest is the ``critical surface'' with $R_c=\infty$, or $m=0$. We
denote it $\Sigma_c$.

Important subspace is the subspace $\Sigma(l)$ generated by applying $RG_l$ to
$\Sigma$ itself,
\[
\Sigma(l)=RG_l\{\Sigma\}\,,
\]
see figure~\ref{fig:nested}.
\begin{figure}[ht]
\centering\figfile{nested}
\caption{The nested images $\Sigma\supset\Sigma(l)\supset\Sigma(l+l')\supset\dots$
of $\Sigma$ under the RG transformations, shrinking to $\Sigma(\infty)$ (thick
curve).}
\label{fig:nested}
\end{figure}
As we have assumed that $RG_l$ acts in $\Sigma$,
\[
\Sigma(l)\subset\Sigma\,.
\]
Obviously the subspaces $\Sigma(l)$ are nested, i.e.
\[
\Sigma(l+l')\subset\Sigma(l)\qquad\text{for any}\qquad l,l'>0
\]
(this follows from $\Sigma(l+l')=RG_l\{\Sigma(l')\}$).

Strictly speaking the space $\Sigma$ is infinite-dimensional. However we assume
that in many important respects it can be treated as finite dimensional one. This
is equivalent to the assumption that it is possible to build a system of
approximations to the RG transformation in which $\Sigma$ is replaced by spaces of
large but finite dimensionalities, and that such approximations converge to exact
RG transformation when this dimensionality goes to infinity.\footnote{Of course
these statements are formulated in a way which is less than satisfactory from
mathematical point of view. Mathematical rigor can be achieved only in some
simplified models of the RG transformation.} If dimensionality of $\Sigma$ is
understood in this sense, the spaces $\Sigma(l)$ for finite $l$ have the same
dimensionality as $\Sigma$ itself. However, the space $\Sigma(\infty)$, obtained
as the limit $l\to\infty$ of $\Sigma(l)$, can be much smaller. The basic idea of
the Wilson's RG approach is in the hope that $\Sigma(\infty)$ has far fewer
dimensions than $\Sigma$. The reason why $\Sigma(\infty)$ is expected to be much
smaller than $\Sigma$ is the following: For large $l$ the action $\mathcal A_l$
describes the physical system with actual interactions of much shorter range than
the range of interactions which appear in the action $\mathcal A_l$ itself.
Remember, $\mathcal A_l$ is physically equivalent to $\mathcal A_0$, with the
scales changed by the factor $e^{-l}$. It is unlikely that arbitrary action in
$\Sigma$ has this property. Thus $\Sigma(\infty)$ is expected to be only a small
subset in $\Sigma$.

The differential equation \eqref{eq:WilsonFlow}, i.e.
\[
\frac{\dd}{\dd l}\,\mathcal A=B\{\mathcal A\}
\]
can be integrated backward in $l$ as well as forward. Suppose $\mathcal A$ is in
$\Sigma$ but not in $\Sigma(l)$. Then $\mathcal A_{-l}$, obtained by integrating
backwards, can not lie in $\Sigma$\footnote{If $\mathcal A_{-l}\in\Sigma$ then
$\mathcal A=RG_l\{\mathcal A_{-l}\}\in\Sigma(l)$.} (figure~\ref{fig:backward}).
\begin{figure}[ht]
\centering\figfile{backward}
\caption{An action $\mathcal A\in\Sigma$ outside $\Sigma(l)$: integrated backwards
by $l$, its trajectory leaves $\Sigma$.}
\label{fig:backward}
\end{figure}

The reason $\mathcal A_{-l}$ is not in $\Sigma$ is likely to be one of the
following: $\mathcal A_{-l}$ can fail to be quasilocal, i.e.\ $\mathcal A_{-l}$
may contain long-range interactions. Or $\mathcal A_{-l}$ can have other
pathological properties, like some coefficients which control convergence of the
functional integral can become such that the functional integral with
$\mathcal A_{-l}$ no longer exists (like $\lambda$ in $\varphi^4$ theory can
become negative), so that the action $\mathcal A_{-l}$ does not define meaningful
theory. Anyhow, for large $l$ the actions $\mathcal A_{-l}$ are typically outside
$\Sigma$.

Important exceptions from this general rule are the actions belonging to the
subspace $\Sigma(\infty)$. If one takes
\[
\mathcal A\in\Sigma(\infty)\,,
\]
and obtains $\mathcal A_{-l}$ by integrating backwards,
\[
\mathcal A_{-l}\in\Sigma(\infty)\,.
\]
This is true, for if $\mathcal A_{-l}$ is in $\Sigma$ but not in
$\Sigma(\infty)$ we have
\[
RG_l\{\mathcal A_{-l}\}=\mathcal A\,,
\]
so that $\mathcal A$ can be reached in finite ``time'' $l$ from some point in
$\Sigma$ outside $\Sigma(\infty)$, but this contradicts the definition of
$\Sigma(\infty)$. In short, starting in $\Sigma(\infty)$ one can integrate back
and forth, still remaining within $\Sigma(\infty)$.

For this reason the subspace $\Sigma(\infty)$ is of central importance in local
Quantum Field Theory. The problem of QFT is to obtain finite local theory in the
limit $\Lambda\to\infty$, with appropriately adjusted ``bare action''
$\mathcal A(\Lambda)$. ``Integrating backward'' is essentially considering the
actions which lead to the same physics as the action $\mathcal A_0$ one starts
with, but equipped with larger and larger cutoff $\Lambda$, measured in the units
of $m$. It is then clear that finding renormalized local theories is equivalent
to finding actions $\mathcal A$ from which one can ``integrate backward'' without
limits. For this reason it is natural to identify $\Sigma(\infty)$ with the space
of local QFT. We will come back to this point many times over.
